Inicialmente la planificación del proyecto plantea cinco partes divididas cada una en tareas con su pertinente estimación temporal como se muestra en la Tabla \ref{tab:estimacion-horas-proyecto}.

\begin{table}[H]
\centering
\footnotesize
\caption{Comparación entre las horas estimadas y las horas reales dedicadas a los bloques y tareas del proyecto.}
\label{tab:estimacion-horas-proyecto}
\renewcommand{\arraystretch}{1.12}
\setlength{\tabcolsep}{5pt}
\begin{tabularx}
{\textwidth}{|>{\raggedright\arraybackslash}X|>{\centering\arraybackslash}p{0.16\textwidth}|>{\centering\arraybackslash}p{0.16\textwidth}|}
\hline
\textbf{Trabajo Fin de Grado} & \cellcolor{yellow!45}\textbf{Horas estimadas} & \cellcolor{yellow!45}\textbf{Horas reales} \\
\hline
\rowcolor{blue!40}\textbf{PT 1 Aprendizaje y desarrollo inicial} & \textbf{62} & \textbf{76} \\
\hline
T 1.1 Análisis del problema e investigación técnica & 24 & 22 \\
\hline
T 1.2 Integración del modelo y de las funciones de pérdida & 22 & 40 \\
\hline
T 1.3 Diseño del preprocesamiento & 16 & 14 \\
\hline
\rowcolor{blue!40}\textbf{PT 2 Entrenamiento} & \textbf{120} & \textbf{144} \\
\hline
T 2.1 Diseño del protocolo de entrenamiento & 10 & 10 \\
\hline
T 2.2 Ajuste de los parámetros de entrenamiento & 28 & 36 \\
\hline
T 2.3 Entrenamiento de dos fuentes y análisis de resultados & 32 & 40 \\
\hline
T 2.4 Entrenamiento de cuatro fuentes y análisis de resultados & 40 & 50 \\
\hline
T 2.5 Organización de los resultados & 10 & 8 \\
\hline
\rowcolor{blue!40}\textbf{PT 3 Implementación de la interfaz} & \textbf{44} & \textbf{39} \\
\hline
T 3.1 Análisis de requisitos & 8 & 6 \\
\hline
T 3.2 Elección de herramientas & 6 & 5 \\
\hline
T 3.3 Diseño del flujo & 14 & 12 \\
\hline
T 3.4 Implementación & 16 & 16 \\
\hline
\rowcolor{blue!40}\textbf{PT 4 Evaluación final} & \textbf{52} & \textbf{87} \\
\hline
T 4.1 Definición del protocolo de evaluación & 10 & 11 \\
\hline
T 4.2 Preparación del conjunto de evaluación & 12 & 12 \\
\hline
T 4.3 Cálculo de métricas & 4 & 4 \\
\hline
T 4.4 Análisis y representación de los resultados & 14 & 14 \\
\hline
T 4.5 Elaboración de conclusiones y trabajo futuro & 12 & 46 \\
\hline
\rowcolor{blue!40}\textbf{PT 5 Documentación} & \textbf{22} & \textbf{64} \\
\hline
T 5.1 Redacción y revisión de la memoria final & 18 & 58 \\
\hline
T 5.2 Limpieza y documentación del código & 4 & 6 \\
\rowcolor{blue!55}\textbf{TIEMPO TOTAL} & \textbf{300} & \textbf{410} \\
\hline
\end{tabularx}
\end{table}

La estimación inicial de tiempo dedicado al proyecto asciende a aproximadamente 300 horas de trabajo efectivo. Esta cifra incluye las tareas de análisis, implementación, depuración, supervisión de entrenamientos, evaluación y documentación. No se contabiliza íntegramente el tiempo durante el cual los modelos se entrenan de forma automática, sino únicamente el esfuerzo dedicado a preparar, supervisar y analizar dichas ejecuciones.
